\section{Incendio o Explosión en Almacén de Material Radiactivo}

\textbf{Clasificación:} 1-A

\subsection*{7.7.1}
Efectúe una verificación con su detector. Asegúrese de portar sus dosímetros de lectura directa y termoluminiscentes.

\subsection*{7.7.2}
Si los niveles de radiación son seguros proceda a inspeccionar en detalle los contenedores de fuentes. Reporte sus resultados al departamento de seguridad radiológica.

\subsection*{7.7.3}
Practique pruebas de fuga a todas las fuentes implicadas.

\subsection*{7.7.4}
Si llegara a encontrar contaminación removible con las pruebas leafrontis, retírese del lugar inspeccionando sus ropas y calzado. Acordone el área y evite la remoción de objetos del área restringida.

\subsection*{7.7.5}
Anote nombres y direcciones de personas que pudieran haber resultado afectadas por la explosión o contaminación e informe a ambulancias y hospitales la presencia de radiación y posible contaminación de los pacientes.

\subsection*{7.7.6}
Evite fumar, ingerir alimentos o bebidas cerca del área del incidente. Notifique a la C.N.S.N.S.

\subsection*{7.7.7}
Si no encontrara contaminación pero los niveles de radiación fueran altos, observe si puede descubrir su procedencia.